\documentclass[11pt]{article}
\usepackage[margin=1in]{geometry}
\usepackage{amsmath,amssymb,amsthm,mathtools}
\usepackage{graphicx}
\usepackage{microtype}
\usepackage[hidelinks]{hyperref}

\newtheorem{theorem}{Theorem}
\newtheorem{lemma}[theorem]{Lemma}
\newcommand{\D}{\mathbb D}
\newcommand{\M}{M}
\newcommand{\Z}{\mathbb Z}
\newcommand{\N}{\mathbb N}
\newcommand{\wh}{\widehat}
\newcommand{\supp}{\operatorname{supp}}

\title{An Infinite Wiener--Pitt Set for the Cantor Group}
\author{Candidate full solution packet\\\small arXiv:1305.3324, Final Remarks item 5}
\date{11 August 2026}

\begin{document}
\maketitle

\noindent\textbf{Status.} Candidate full affirmative answer, subject to human
review.  The proof uses the finite-group integer-rounding theorem of
Ohrysko--Sanders--Wojciechowski \cite[Theorem 11]{OSW}; the passage from finite
quotients to the Cantor group and the norm-diagonal construction below are the
new steps.

\section{Source question}

Ohrysko and Wojciechowski \cite[Final Remarks, item 5]{OW} ask whether any
infinite Wiener--Pitt set exists for the convolution measure algebra of the
Cantor group.  Their printed question appears on journal page 25:

\begin{center}
\includegraphics[width=\textwidth]{figures/open_problem_crop.png}
\end{center}

We use the notation
\[
  \D:=\prod_{r\geq1}\Z/2\Z,
  \qquad \Gamma:=\wh\D=\bigoplus_{r\geq1}\Z/2\Z.
\]
For $\mu\in \M(\D)$, its Fourier--Stieltjes transform is
$\wh\mu(\gamma)=\int_{\D}\overline{\gamma(x)}\,d\mu(x)$.
A compact set $K\subset\mathbb C$ is a \emph{Wiener--Pitt set} for
$\M(\D)$ if
\[
 \wh\mu(\Gamma)\subset K
 \quad\Longrightarrow\quad
 \sigma_{\M(\D)}(\mu)=\overline{\wh\mu(\Gamma)}
\]
for every $\mu\in\M(\D)$.

\section{Proof intuition}

The obstruction identified in the source paper is that one Fourier value can
occur on an enormous set of characters.  On the Cantor group this is not a
reason to count the level set: the dual group is the increasing union of
finite groups.  We round on every finite subgroup, lift the resulting measure
from the corresponding finite quotient, and take a weak-star limit.  This
turns every sufficiently separated Fourier level into an idempotent measure,
with the same quantitative norm bound as in the finite theorem.

A second obstruction is that the definition of a Wiener--Pitt set places no
bound on $\|\mu\|$.  We handle all norms diagonally.  At stage $N$ an
interpolation polynomial labels the first $N$ Fourier levels by the integers
$1,\ldots,N$ and sends a sufficiently tiny tail close to zero.  For measures
of norm at most $N$, rounding this one polynomial isolates the whole prefix
and leaves a residual small enough to start an infinite idempotent expansion.

The expansion of $\mu$ itself need not converge in measure norm.  Its square
and cube do.  A character of the measure algebra can select at most one of the
orthogonal level idempotents.  The square determines a level up to sign, and
the cube fixes the sign, forcing every spectral value to be an actual Fourier
value (or a limit of such values).

\section{A rounding lemma for the Cantor group}

For a real-valued function $f$ with
$d(f,\Z):=\sup_x\min_{m\in\Z}|f(x)-m|<1/2$, write $f_{\Z}$ for its
pointwise nearest-integer function.  We use the following form of
\cite[Theorem 11]{OSW}: there is a function
\[
 \delta:(0,1]\times(0,1]\times[1,\infty)\longrightarrow(0,1/2)
\]
such that, for every finite abelian group $F$, every real
$f\in A(F)$ satisfying $\|f\|_{A(F)}\leq M$ and
$d(f,\Z)\leq\delta(\varepsilon_1,\varepsilon_2,M)$ obeys
\[
 \|f_{\Z}\|_{A(F)}
 \leq(1+\varepsilon_1)\|f\|_{A(F)}+\varepsilon_2.
\]

\begin{lemma}[profinite integer rounding]\label{lem:round}
Let $\theta\in\M(\D)$ have real Fourier--Stieltjes transform and
$\|\theta\|\leq M$, where $M\geq1$.  If
\[
 d(\wh\theta,\Z)\leq
 \delta(\varepsilon_1,\varepsilon_2,M),
\]
then there is $\theta_{\Z}\in\M(\D)$ such that
\[
 \wh{\theta_{\Z}}=(\wh\theta)_{\Z},
 \qquad
 \|\theta_{\Z}\|
 \leq(1+\varepsilon_1)\|\theta\|+\varepsilon_2.
\]
No finiteness assumption is made on the support of
$(\wh\theta)_{\Z}$.
\end{lemma}

\begin{proof}
Let $\Gamma_r\leq\Gamma$ be the subgroup generated by the first $r$
coordinates, let $H_r:=\Gamma_r^\perp$, and let
$\pi_r:\D\to\D/H_r$ be the quotient map.  The finite measure
$(\pi_r)_\#\theta$ has norm at most $\|\theta\|$, and its Fourier
transform is $\wh\theta|_{\Gamma_r}$.  Applying the finite-group theorem
to this restriction gives a measure $\lambda_r$ on $\D/H_r$ whose Fourier
transform is $(\wh\theta|_{\Gamma_r})_{\Z}$ and such that
\[
 \|\lambda_r\|
 \leq(1+\varepsilon_1)\|\theta\|+\varepsilon_2.
\]

Lift $\lambda_r$ to a measure $\widetilde\lambda_r$ on $\D$ by replacing
each atom at a coset $x+H_r$ with the same scalar multiple of normalized Haar
measure on that coset.  Then $\|\widetilde\lambda_r\|=\|\lambda_r\|$ and
\[
 \wh{\widetilde\lambda_r}(\gamma)=
 \begin{cases}
   (\wh\theta(\gamma))_{\Z},&\gamma\in\Gamma_r,\\
   0,&\gamma\notin\Gamma_r.
 \end{cases}
\]
The uniformly bounded net $(\widetilde\lambda_r)$ has a weak-star cluster
point $\theta_{\Z}$ in $\M(\D)=C(\D)^*$.  For each fixed
$\gamma\in\Gamma$, the displayed Fourier coefficient is equal to
$(\wh\theta(\gamma))_{\Z}$ for all sufficiently large $r$.  Hence the
cluster point has the required transform, and weak-star lower
semicontinuity gives the norm bound.
\end{proof}

\section{Construction of the compact set}

We recursively construct a strictly decreasing sequence
\[
 1=a_1>a_2>\cdots>0
\]
and put $K:=\{0\}\cup\{a_n:n\geq1\}$.

Suppose $a_1,\ldots,a_N$ have been chosen.  Let $q_N$ be the unique real
polynomial of degree at most $N$ satisfying
\[
 q_N(0)=0,
 \qquad q_N(a_j)=j\quad(1\leq j\leq N).
\]
Write $q_N(t)=\sum_k c_{N,k}t^k$ and set
\[
 B_N:=\max\left\{1,\sum_k|c_{N,k}|N^k\right\},
 \qquad E_N:=2B_N+1.
\]
For $1\leq j\leq N$, let $p_{j,N}$ be the interpolation polynomial on
$\{0,1,\ldots,N\}$ satisfying
\[
 p_{j,N}(j)=1,
 \qquad p_{j,N}(k)=0\quad(k\in\{0,1,\ldots,N\},\ k\ne j).
\]
If $p_{j,N}(t)=\sum_k d_{j,N,k}t^k$, define
\[
 C_N:=\max_{1\leq j\leq N}\sum_k|d_{j,N,k}|E_N^k,
 \qquad
 T_N:=N+C_N\sum_{j=1}^N a_j.
\]
All these numbers are determined before $a_{N+1}$ is chosen.  Since
$q_N(0)=0$, we may choose $a_{N+1}>0$ so small that
\begin{align}
 a_{N+1}&<a_N/9,\label{eq:ratio}\\
 \frac{a_{N+1}}{a_N}
 &\leq\delta\!\left(\tfrac12,\tfrac12,a_N^{-3/2}\right),
 \label{eq:tailround}\\
 \sup_{0\leq t\leq a_{N+1}}|q_N(t)|
 &\leq\delta(1,1,B_N),\label{eq:prefixround}\\
 a_{N+1}&\leq T_N^{-2}.\label{eq:start}
\end{align}
Every right-hand side is positive, so the recursion is possible.  Condition
\eqref{eq:ratio} implies $a_n\to0$; hence $K$ is compact and infinite.

\section{Main theorem}

\begin{theorem}\label{thm:main}
The set $K$ constructed above is a Wiener--Pitt set for $\M(\D)$.  In
particular, the question in \cite[Final Remarks, item 5]{OW} has an
affirmative answer.
\end{theorem}

\begin{proof}
Let $\mu\in\M(\D)$ satisfy $\wh\mu(\Gamma)\subset K$, and choose an integer
$N\geq\max\{1,\|\mu\|\}$.  Put
\[
 A_j:=\{\gamma\in\Gamma:\wh\mu(\gamma)=a_j\}.
\]
Because $K\subset\mathbb R$, the transform of $\mu$ is real.

\subsection*{The finite prefix}
The polynomial estimate gives $\|q_N(\mu)\|\leq B_N$.  By
\eqref{eq:prefixround}, $q_N(\wh\mu)$ is within
$\delta(1,1,B_N)$ of the integers: it equals $j$ on $A_j$ for
$j\leq N$, and has absolute value at most that threshold on every later
level.  Lemma~\ref{lem:round} therefore gives $\eta_N\in\M(\D)$ with
\[
 \wh\eta_N=\sum_{j=1}^N j\,1_{A_j},
 \qquad \|\eta_N\|\leq2B_N+1=E_N.
\]
Set
\[
 e_j:=p_{j,N}(\eta_N)\qquad(1\leq j\leq N).
\]
Then $\wh e_j=1_{A_j}$ and $\|e_j\|\leq C_N$.  In particular, the
$e_j$ are idempotent and pairwise convolution-orthogonal.  The residual
\[
 \tau_{N+1}:=\mu-\sum_{j=1}^N a_je_j
\]
satisfies
\[
 \|\tau_{N+1}\|
 \leq\|\mu\|+C_N\sum_{j=1}^Na_j
 \leq T_N\leq a_{N+1}^{-1/2},
\]
and its transform is $a_j1_{A_j}$ on levels $j>N$ and zero elsewhere.

\subsection*{The infinite tail}
Inductively suppose $n\geq N+1$, the idempotents $e_j$ have been defined
for $j<n$, and
\[
 \tau_n:=\mu-\sum_{j<n}a_je_j,
 \qquad \|\tau_n\|\leq a_n^{-1/2}.
\]
The transform of $a_n^{-1}\tau_n$ is $1$ on $A_n$, is zero on the
earlier levels, and has absolute value at most $a_{n+1}/a_n$ on all later
levels.  Its norm is at most $a_n^{-3/2}$.  By
\eqref{eq:tailround} and Lemma~\ref{lem:round} with
$\varepsilon_1=\varepsilon_2=1/2$, there is an idempotent $e_n$ with
\[
 \wh e_n=1_{A_n},
 \qquad
 \|e_n\|\leq\tfrac32a_n^{-3/2}+\tfrac12.
\]
Define $\tau_{n+1}:=\tau_n-a_ne_n$.  Since $a_n\leq1$,
\[
 \|\tau_{n+1}\|
 \leq a_n^{-1/2}+a_n\!\left(\tfrac32a_n^{-3/2}+\tfrac12\right)
 \leq3a_n^{-1/2}
 \leq a_{n+1}^{-1/2},
\]
where the last inequality follows from \eqref{eq:ratio}.  This completes
the induction.  Fourier uniqueness also shows that all $e_j$ are pairwise
orthogonal, including across the prefix/tail boundary.

\subsection*{Norm-convergent powers}
For $m=2,3$, the series $\sum_j a_j^m e_j$ converges in measure norm.
Indeed, the prefix is finite and, on the tail,
\[
 a_j^m\|e_j\|
 \leq\tfrac32a_j^{m-3/2}+\tfrac12a_j^m,
\]
whose sum converges because $a_{j+1}<a_j/9$.  Its Fourier transform is
$\wh\mu^m$, so Fourier uniqueness yields
\begin{equation}\label{eq:powers}
 \mu^{*m}=\sum_{j\geq1}a_j^m e_j,
 \qquad m=2,3.
\end{equation}

Let $\varphi$ be a character of the commutative Banach algebra $\M(\D)$
and put $z:=\varphi(\mu)$.  Each $\varphi(e_j)$ is $0$ or $1$, and
orthogonality permits at most one index with value $1$.  If all are zero,
\eqref{eq:powers} gives $z^2=0$, hence $z=0$.  Otherwise, for the unique
index $j$ with $\varphi(e_j)=1$, it gives
\[
 z^2=a_j^2,
 \qquad z^3=a_j^3,
\]
and therefore $z=a_j$.  Moreover, $\varphi(e_j)=1$ implies $e_j\ne0$;
as $\wh e_j=1_{A_j}$, Fourier uniqueness then implies $A_j\ne\varnothing$.
Thus every nonzero spectral value is an actual Fourier value.

It remains only to handle a possible spectral value $0$.  If $0$ lies in
$\overline{\wh\mu(\Gamma)}$, there is nothing to prove.  Otherwise only
finitely many levels $A_j$ are nonempty and $0$ is not itself a Fourier
value.  If $J$ is that finite set of indices, the polynomial
$P(t)=\prod_{j\in J}(t-a_j)$ satisfies $\wh{P(\mu)}=0$, hence
$P(\mu)=0$.  Spectral mapping then excludes $0$ from $\sigma(\mu)$ because
$P(0)\ne0$.  Consequently
\[
 \sigma(\mu)\subset\overline{\wh\mu(\Gamma)}.
\]
The reverse inclusion always holds because Fourier evaluations are
characters and the spectrum is closed.  Hence equality holds, completing
the proof.
\end{proof}

\section{Verification and novelty status}

The proof is qualitative: it uses only positivity of the rounding thresholds,
so no numerical estimate of $a_n$ is required.  The recursive order is
non-circular: $q_N,B_N,E_N,C_N,T_N$ depend only on
$a_1,\ldots,a_N$, and all four conditions on $a_{N+1}$ can therefore be
imposed simultaneously.

The two principal audit points are:
\begin{enumerate}
\item the lifted finite-quotient measures have exactly zero Fourier transform
outside $\Gamma_r$, and weak-star convergence fixes every individual
coefficient eventually;
\item the cube in \eqref{eq:powers} is essential: the square alone would allow
the extraneous value $-a_j$.
\end{enumerate}

A bounded novelty search was performed through 11 August 2026 in the run
indexes, arXiv, the journal/publisher record, and exact web searches using
``Cantor group'', ``infinite Wiener--Pitt set'', arXiv:1305.3324, the source
authors, and arXiv:2510.24578.  The only direct continuation found is
\cite{OSW}, which assumes both strong continuity and $\|\mu\|\leq1$ and
explicitly says its method does not remove strong continuity.  No all-measures
answer to the Cantor-group question was found.  The novelty claim remains
provisional pending expert literature review.

\section{Human-review recommendation}

Prioritize review of Lemma~\ref{lem:round}, especially the identification of
the finite-group $A$-norm with the total variation norm of the pushed-forward
measure, and the simultaneous prefix isolation at stage $N$.  If those steps
are accepted, the remaining induction and Gelfand-spectrum argument are
elementary and close the question as stated.

\begin{thebibliography}{9}
\bibitem{OW}
P. Ohrysko and M. Wojciechowski,
\emph{On the relationships between Fourier--Stieltjes coefficients and spectra
of measures}, Studia Math. 221 (2014), 117--140; arXiv:1305.3324.

\bibitem{OSW}
P. Ohrysko, T. Sanders, and M. Wojciechowski,
\emph{Wiener--Pitt sets for compact Abelian groups}, J. Fourier Anal. Appl.
32 (2026), Article 30; arXiv:2510.24578; DOI
\href{https://doi.org/10.1007/s00041-026-10239-1}{10.1007/s00041-026-10239-1}.
\end{thebibliography}

\end{document}
